\documentclass[10pt]{article}
\usepackage[margin=0.75in]{geometry}
\usepackage{listings}
\usepackage{xcolor}
\usepackage{titlesec}
\setlength{\parskip}{2pt}
\setlength{\parindent}{0pt}
\titlespacing*{\section}{0pt}{4pt}{2pt}
\titleformat*{\section}{\normalsize\bfseries}

\lstset{
    basicstyle=\ttfamily\scriptsize,
    breaklines=true,
    columns=fullflexible,
    keywordstyle=\color{blue},
    commentstyle=\color{gray},
    stringstyle=\color{red!70!black},
    showstringspaces=false,
    language=Python,
    frame=single,
    framerule=0.3pt,
    rulecolor=\color{black!30},
    xleftmargin=0.5em,
    xrightmargin=0.5em,
    aboveskip=2pt,
    belowskip=2pt
}

\begin{document}
\begin{center}
{\large\bfseries Homework 3: HL7 Parsing and Modification}\\
Jonathan DeLemos \quad (PennKey: delemos1)
\end{center}
\vspace{-0.5em}

\section*{Summary of Approach}
The instructions for this assignment were clear and well-scoped. After reading them through, I completed the code in the order the template laid out: read \texttt{sample-hl7.txt}, normalize line endings, parse with the \texttt{hl7} Python library, extract the requested fields into a dictionary, dump it to \texttt{output.json}, then mutate the in-memory message to apply the required changes and write the result to \texttt{modified-hl7.txt}. I used \texttt{hl7inspector.com} to confirm the field indices for each segment (PID, PV1, OBR, DG1, OBX) before writing the access code.

\section*{HL7 vs.\ Other Text-Based Data, and vs.\ DICOM}
HL7 looks like text but behaves more like a delimited, positionally-encoded record format. Unlike JSON, CSV, or XML, there are no field names in the payload --- every field is identified by its ordinal position within a pipe-delimited segment, and the meaning of each position is defined externally by the HL7 standard and message version. There is also a five-level nesting hierarchy (segment $\rightarrow$ field $\rightarrow$ repetition $\rightarrow$ component $\rightarrow$ sub-component) using several different delimiters (\verb|\r|, \verb!|!, \verb|~|, \verb|^|, \verb|&|), and segments terminate on carriage returns rather than newlines, which is unusual on modern systems. DICOM, by contrast, is a binary, tag-based format: each element is identified by a numeric \texttt{(group, element)} tag with an explicit value representation describing its type, length, and encoding, and the payload can carry pixel data alongside metadata. Where HL7 is a structured text telegram designed for messaging between information systems, DICOM is a self-describing binary container designed for imaging and persistence --- so DICOM effectively requires a library to read, while HL7 \emph{could} be parsed with string operations (though doing so is unwise given the delimiter hierarchy and escape rules).

\section*{Accommodations in My Code}
Two accommodations were needed. First, the source file used lone carriage returns as segment terminators, so the parser would not split segments correctly until I explicitly replaced \verb|\n| with \verb|\r| in the loaded string (Python's universal-newlines mode converts \verb|\r| to \verb|\n| on read). Second, several fields are positional but only meaningful in a deeper component --- e.g., OBR-31 stores the reason for study in the \emph{second} component of the field (the first is empty), so I had to drill in with \texttt{obr[31][0][1]} rather than just \texttt{obr[31]}. For the patient name I extracted the family and given names from PID-5 as separate components and assembled them into a human-readable ``First Last'' string, while leaving the referring doctor field as-is per the assignment instructions.

\section*{Highest Source of Errors}
The biggest hindrance was unfamiliarity with the \texttt{hl7} Python package --- specifically, getting comfortable with the shape of the outputs returned by its accessors. The library returns nested container types (\texttt{Message}, \texttt{Segment}, \texttt{Field}, \texttt{Repetition}, \texttt{Component}) rather than plain strings, so my early attempts produced values that looked correct when printed but were actually wrapped objects with subtly wrong sub-indexing. For example, my first attempt at the reason-for-study extraction was:

\begin{lstlisting}
reason_for_study = obr[31]   # returns a Field, not a string
\end{lstlisting}

\noindent which serialized as \texttt{\^{}Abd pain and weight loss, non-acute} --- including the leading caret because the first component was empty. The text actually lives in the second component of the first repetition, so the correct access is:

\begin{lstlisting}
reason_for_study = str(obr[31][0][1])   # 'Abd pain and weight loss, non-acute'
\end{lstlisting}

\noindent I resolved this class of bug by writing a small exploration script that printed \texttt{repr()} of each field I needed, which revealed the exact nesting depth (e.g.\ \texttt{PID[5]} surfaced as \texttt{[[['TIGER'], ['DANIELLE']]]}, telling me I needed two index steps to reach the family name and to wrap the result in \texttt{str()} for clean JSON output). Once I had that mental model of the package's return types, the rest of the extractions and in-place modifications fell out cleanly.

\end{document}
